\section{Metric, coordinates, and boundary layers}
\label{sec:mathematics}

\subsection{Face-turn metric and subgroup coordinates}
\label{sec:metric-coordinates}

Let a face symbol denote a clockwise turn through one fifth of a revolution.
For a set $F$ of faces, write
\[
  \Sigma_F=\{X^k:X\in F,\ 1\leq k\leq4\}.
\]
Every element of $\Sigma_F$ has cost one.  Thus $X$, $X^2$, $X^3$, and
$X^4=X^{-1}$ are each one move in the face-turn metric (FTM); rotations,
slice turns, and identity moves are not in the alphabet.  In particular,
$\Sigma_F$ is symmetric.  If $w=s_1\cdots s_m$ is a word, then
$|w|=m$, and its tokens act from left to right:
\[
  x\mathbin{\cdot}w
  =(((x\mathbin{\cdot}s_1)\mathbin{\cdot}s_2)\cdots)
    \mathbin{\cdot}s_m.
\]
This convention matters for the physical lifting below.

Let $\Omega$ denote the full $48$-move FTM alphabet obtained by taking all
four nonidentity powers of all $12$ faces.  Words over $\Omega$ may temporarily
leave a subgroup; only their start and target memberships matter for an
ambient FTM bound.

The part of the subgroup chain used in this work is
\begin{align*}
  G_5&=\langle U,R,F,L,\mathrm{BR},\mathrm{BL},\mathrm{FR}\rangle,
  &G_6&=\langle U,R,F,L,\mathrm{BR},\mathrm{BL}\rangle,\\
  G_7&=\langle U,R,F,L,\mathrm{BR}\rangle,
  &G_8&=\langle U,R,F,L\rangle,\\
  G_9&=\langle U,R,F\rangle.
\end{align*}
For $G_i$ we use the FTM alphabet containing the four nonidentity powers of
each displayed generator of $G_i$; denote it by $\Sigma_i$.  Hence
$\Sigma_{i+1}\subseteq\Sigma_i$ for the transitions considered here.

\begin{definition}[Target-solving distance]
\label{def:target-distance}
Let $A\geq B$ be two groups in the chain, and let $\Gamma$ be a symmetric FTM
alphabet.  For a state $x\in A$, define
\[
  d_{\Gamma}(x,B)
    =\min\bigl\{|w|:w\in\Gamma^*,\ x\mathbin{\cdot}w\in B\bigr\},
\]
and define the worst-case target-solving radius over starts in $A$ by
\[
  \Rsolve_{\Gamma}(A,B)=\max_{x\in A}d_{\Gamma}(x,B).
\]
For a phase table we use $\Gamma=\Sigma_A$ and abbreviate
\begin{align*}
  \rho(A,B)&=\Rsolve_{\Sigma_A}(A,B),\\
  d_{A,B}(Bx)&=d_{\Sigma_A}(x,B),\\
  \mathcal M(A,B)
    &=\{Bx\in B\backslash A:d_{A,B}(Bx)=\rho(A,B)\}.
\end{align*}
\end{definition}

The distance is independent of the representative of a left $B$-coset: if
$x'=bx$ with $b\in B$, then $x'\mathbin{\cdot}w\in B$ if and only if
$x\mathbin{\cdot}w\in B$.  When $\Gamma=\Sigma_A$, $\rho(A,B)$ is the
eccentricity of the target vertex $B$ in the Schreier graph on
$B\backslash A$ generated by $\Sigma_A$.  It is not, by definition, the
ordinary graph diameter, which maximizes distance over arbitrary pairs of
vertices.  We therefore use ``target-solving radius'' rather than
``diameter'' for the quantities proved by the phase tables.  The combined
claims below allow $\Gamma=\Sigma_5$, which may be larger than the source
alphabet, and also state the resulting ambient $\Gamma=\Omega$ bounds.  Their
solution paths need not remain inside the source subgroup and hence should not
be described as diameters of those phase Schreier graphs.

\subsection{A complete physical lift of two coordinates}
\label{sec:physical-lift}

Let $A\geq B\geq C$.  For each $p\in B\backslash A$, fix an optimal lowering
word $u_p\in\Sigma_A^*$ such that
\[
  p\mathbin{\cdot}u_p=B,
  \qquad |u_p|=d_{A,B}(p).
\]
The first equality implies $p=B u_p^{-1}$.  For each
$q\in C\backslash B$, fix a physical representative $r_q\in B$ with
$Cr_q=q$.  Define the lifted state
\begin{equation}
  \lambda(p,q)=r_q\mathbin{\cdot}u_p^{-1}.
  \label{eq:physical-lift}
\end{equation}
Thus, under left-to-right replay, the product order is the representative of
the \emph{second} coordinate followed by the inverse of the lowering word for
the \emph{first} coordinate.

\begin{lemma}[Lifting and product order]
\label{lem:maximal-lifting}
For every $(p,q)\in(B\backslash A)\times(C\backslash B)$, the state
$\lambda(p,q)$ satisfies
\begin{align}
  B\lambda(p,q)&=p,
  \label{eq:lift-first-coordinate}\\
  \lambda(p,q)\mathbin{\cdot}u_p&=r_q,
  \label{eq:lift-intermediate}\\
  C\bigl(\lambda(p,q)\mathbin{\cdot}u_p\bigr)&=q.
  \label{eq:lift-second-coordinate}
\end{align}
Moreover, the lift is complete in the following sense.  Given any $x\in A$,
put $p=Bx$ and $q=C(x\mathbin{\cdot}u_p)$.  Then
\begin{equation}
  Cx=C\lambda(p,q).
  \label{eq:lift-completeness}
\end{equation}
Consequently, one lifted state for each pair $(p,q)$ covers every combined
$C$-coset that can arise after the chosen optimal first-phase lowering word.
Injectivity is neither assumed nor required.
\end{lemma}

\begin{proof}
Because $r_q\in B$ and $p=B u_p^{-1}$,
\[
  B\lambda(p,q)=B r_q u_p^{-1}=B u_p^{-1}=p.
\]
Applying $u_p$ cancels its inverse in~\eqref{eq:physical-lift}, giving
$r_q$ and hence the second coordinate $Cr_q=q$.

For completeness, $p=Bx$ and $p\mathbin{\cdot}u_p=B$ imply that
$y=x\mathbin{\cdot}u_p$ belongs to $B$.  If $q=Cy$, then $y$ and $r_q$
are in the same left $C$-coset, so $y=cr_q$ for some $c\in C$.  Therefore
\[
  x=y\mathbin{\cdot}u_p^{-1}
    =c r_q\mathbin{\cdot}u_p^{-1}
    =c\lambda(p,q),
\]
which proves~\eqref{eq:lift-completeness}.
\end{proof}

The artifact materializes~\eqref{eq:physical-lift} in a full physical
Megaminx state rather than combining coordinate integers.  It then checks
\cref{eq:lift-first-coordinate,eq:lift-intermediate,eq:lift-second-coordinate}
for every raw pair.  Although the proof does not need injectivity, the two
concrete lifted products in this work happen to contain no duplicate physical
states.

\subsection{Boundary-layer reduction}
\label{sec:boundary-layer-reduction}

\begin{lemma}[Boundary-layer reduction]
\label{lem:boundary-layer}
Let $A\geq B\geq C$ be finite groups with symmetric generating alphabets
$\Sigma_B\subseteq\Sigma_A$.  Let $\Gamma$ be any symmetric FTM alphabet
containing $\Sigma_A$, and put
\[
  a=\rho(A,B),
  \qquad b=\rho(B,C).
\]
Construct $\lambda(p,q)$ as in~\cref{lem:maximal-lifting}.  Suppose that
for every
\[
  (p,q)\in\mathcal M(A,B)\times\mathcal M(B,C)
\]
there is a word over $\Gamma$ of length at most $a+b-1$ that takes
$\lambda(p,q)$ into $C$.  Then
\[
  \Rsolve_{\Gamma}(A,C)\leq a+b-1.
\]
\end{lemma}

\begin{proof}
Fix $x\in A$ and set $p=Bx$.  If $d_{A,B}(p)\leq a-1$, first apply the
optimal word $u_p$ and then solve the resulting element of $B$ into $C$ in at
most $b$ moves.  The total is at most $a+b-1$.

It remains to consider $p\in\mathcal M(A,B)$, so $|u_p|=a$.  Put
$y=x\mathbin{\cdot}u_p\in B$ and $q=Cy$.  If
$d_{B,C}(q)\leq b-1$, the same sequential argument costs at most
$a+b-1$.  Otherwise $q\in\mathcal M(B,C)$.  By
\cref{lem:maximal-lifting}, $Cx=C\lambda(p,q)$.  Ambient target-solving
distance to $C$ depends only on this left coset, and the assumed frontier word
therefore solves $x$ in at most $a+b-1$ moves.  These cases exhaust $A$.
\end{proof}

The lemma isolates precisely the simultaneous-maximal boundary.  It does not
claim that a search over this boundary is complete: completeness comes from
enumerating the entire product and replaying a valid bounded word for each
member.

\subsection{Exact phase data and exhaustive frontiers}
\label{sec:phase-data}

The proof treats the implementation coordinates as group quotients only after
an exact coordinate-contract audit.  For each phase, the coordinate is the
orbit of a tuple of labelled, oriented pieces under the source generators.
The builder checks that every target generator fixes the target tuple, that
the coordinate ranks and transitions are well defined, and that the BFS
reaches the full declared coordinate set.  A Schreier--Sims audit on the
pinned $120$-position actions recomputes
\begin{align*}
 [G_5:G_7]&=208{,}099{,}584\cdot68{,}400,\\
 [G_7:G_9]&=64{,}157{,}184\cdot25{,}945{,}920,
\end{align*}
as well as the two intermediate subgroup orders.  To see why these checks
suffice, let $n_1,n_2$ be the two coordinate counts in either pair.  Target
invariance makes each coordinate orbit a quotient by a stabilizer containing
the stated target subgroup, so
\(n_1\leq[A:B]\) and \(n_2\leq[B:C]\).  The exact order audit gives
\(n_1n_2=[A:C]=[A:B][B:C]\); hence both inequalities are equalities and the
target stabilizers are exactly $B$ and $C$.  Thus the coordinate fibers are
the left cosets used in \cref{def:target-distance}, rather than merely hashes
that agree in sampled tests.

\begin{proposition}[Phase-table data]
\label{prop:phase-data}
Complete breadth-first phase tables, regenerated from the pinned move and
coordinate definitions, establish the following target-solving data.
\begin{center}
\begingroup
\small
\setlength{\tabcolsep}{3.5pt}
\begin{tabularx}{\textwidth}{@{}cc>{\raggedright\arraybackslash}Xrrrr@{}}
\toprule
Phase & Transition & Source faces & $|\Sigma_A|$ & $|B\backslash A|$
  & $\rho(A,B)$ & $|\mathcal M(A,B)|$ \\
\midrule
3 & $G_5\to G_6$ & $U,R,F,L,\mathrm{BR},\mathrm{BL},\mathrm{FR}$
  & 28 & $208{,}099{,}584$ & 14 & 212 \\
4 & $G_6\to G_7$ & $U,R,F,L,\mathrm{BR},\mathrm{BL}$
  & 24 & $68{,}400$ & 8 & $2{,}531$ \\
5 & $G_7\to G_8$ & $U,R,F,L,\mathrm{BR}$
  & 20 & $64{,}157{,}184$ & 13 & $3{,}484$ \\
6 & $G_8\to G_9$ & $U,R,F,L$
  & 16 & $25{,}945{,}920$ & 13 & 117 \\
\bottomrule
\end{tabularx}
\endgroup
\end{center}
Here the radii are distances from the target cosets, in the sense of
\cref{def:target-distance}; no ordinary Schreier-graph diameter is asserted.
\end{proposition}

For phases 3 and 4, the combined quotient has
\begin{equation}
  |G_7\backslash G_5|
    =208{,}099{,}584\cdot68{,}400
    =14{,}234{,}011{,}545{,}600,
  \label{eq:pair34-quotient-size}
\end{equation}
whereas the simultaneous-maximal product contains only
\begin{equation}
  212\cdot2{,}531=536{,}572
  \label{eq:pair34-frontier-size}
\end{equation}
raw lifts.  For phases 5 and 6, the corresponding values are
\begin{align}
  |G_9\backslash G_7|
    &=64{,}157{,}184\cdot25{,}945{,}920
      =1{,}664{,}617{,}163{,}489{,}280,
  \label{eq:pair56-quotient-size}\\
  3{,}484\cdot117&=407{,}628.
  \label{eq:pair56-frontier-size}
\end{align}
The complete physical lifting check found respectively $536{,}572$ and
$407{,}628$ valid, pairwise distinct full states (maximum multiplicity one).
The absence of duplicates makes the displayed products the exact numbers of
stored obligations; the logical coverage in~\cref{lem:maximal-lifting} would
remain valid even if duplicates had occurred.

\begin{proposition}[Exhaustive certificate coverage]
\label{prop:certificate-coverage}
The regenerated maximal-layer products admit the following exact partitions.
\begin{center}
\begin{tabular}{@{}lrrrr@{}}
\toprule
Pair & Raw lifts & Exact-rewrite words & Direct words & Missing \\
\midrule
$3+4$ & $536{,}572$ & 203 & $536{,}369$ & 0 \\
$5+6$ & $407{,}628$ & $8{,}461$ & $399{,}167$ & 0 \\
\bottomrule
\end{tabular}
\end{center}
Across the two pairs this is an exhaustive partition
\[
  944{,}200=8{,}664+935{,}536
\]
into exact-rewrite and direct-certificate obligations.
For pair $3+4$, every replayed word has length at most 21 and in fact uses the
restricted alphabet $\Sigma_5$.  For pair $5+6$, every replayed word has
length at most 25: the $8{,}461$ exact-rewrite witnesses are words over
$\Sigma_5$, while all $399{,}167$ direct words use the restricted alphabet
$\Sigma_7\subseteq\Sigma_5$.  For every ID, the checker reconstructs the
prescribed physical lift and requires membership in exactly one part of the
partition.  Exact-rewrite words are replayed in the full physical
representation all the way to the solved state.  Direct words are replayed
independently in the combined quotient and in the full physical
representation to the appropriate target subgroup.  Thus the proposition
depends on the words and their exhaustive ID coverage, not on the method that
discovered them.
\end{proposition}

\begin{remark}[Alphabet scope]
\label{rem:alphabet-scope}
The distinction between $\Sigma_7$ and $\Sigma_5$ is substantive.  The neural
database alone covers only the complement of the exact-rewrite set, so its
$\Sigma_7$ words do not establish a restricted-alphabet bound for every
frontier state.  Pair $5+6$ preprocessing uses the first seven faces, so the
complete partition establishes a $\Sigma_5$ bound, but not a $\Sigma_7$
bound.  The ambient conclusion follows because
$\Sigma_7\subseteq\Sigma_5\subseteq\Omega$.  Pair $3+4$ likewise has a
$\Sigma_5$ conclusion because both parts of its partition use that alphabet.
\end{remark}

\subsection{Certified upper bounds}
\label{sec:certified-bounds}

\begin{theorem}[Phases 3 and 4]
\label{thm:pair34}
The combined target-solving radii satisfy
\[
  \Rsolve_{\Sigma_5}(G_5,G_7)\leq21,
  \qquad
  \Rsolve_{\Omega}(G_5,G_7)\leq21.
\]
\end{theorem}

\begin{proof}
By~\cref{prop:phase-data}, the consecutive radii are 14 and 8.  The
simultaneous-maximal product is exhausted by
\cref{eq:pair34-frontier-size}, and every member has a word of length at most
$21=14+8-1$ by~\cref{prop:certificate-coverage}.  Apply
\cref{lem:boundary-layer} with $(A,B,C)=(G_5,G_6,G_7)$ and
$\Gamma=\Sigma_5$.  The ambient statement follows because
$\Sigma_5\subseteq\Omega$.
\end{proof}

\begin{theorem}[Phases 5 and 6]
\label{thm:pair56}
The combined target-solving radii satisfy
\[
  \Rsolve_{\Sigma_5}(G_7,G_9)\leq25,
  \qquad
  \Rsolve_{\Omega}(G_7,G_9)\leq25.
\]
\end{theorem}

\begin{proof}
The two consecutive radii are 13 and 13 by~\cref{prop:phase-data}.  The
product in~\cref{eq:pair56-frontier-size} is exhaustive, and
\cref{prop:certificate-coverage} supplies a word of length at most
$25=13+13-1$ for every member.  The claim follows from
\cref{lem:boundary-layer} with $(A,B,C)=(G_7,G_8,G_9)$ and
$\Gamma=\Sigma_5$.  The ambient statement follows because
$\Sigma_5\subseteq\Omega$.
\end{proof}

Theorems~\ref{thm:pair34} and~\ref{thm:pair56} are upper bounds only.  We do
not provide lower-bound witnesses at distances 21 or 25, and therefore do not
claim exact combined radii or ordinary graph diameters.

\begin{corollary}[Conditional 112-move bound]
\label{cor:conditional-112}
Assume that the externally reported 114-move Megaminx phase chain is a valid
FTM upper bound under the same move convention, subgroup chain, and phase
composition semantics used here, and that it contains the separate
contributions $14+8$ for phases 3 and 4 and $13+13$ for phases 5 and 6.  Then
every Megaminx state is solvable in at most 112 FTM moves.
\end{corollary}

\begin{proof}
Replacing $14+8=22$ by~\cref{thm:pair34}'s bound 21 saves one move, and
replacing $13+13=26$ by~\cref{thm:pair56}'s bound 25 saves a second.  All
other contributions in the assumed chain are unchanged, so its total becomes
$114-1-1=112$.
\end{proof}

\begin{corollary}[Conditional 111-move bound]
\label{cor:conditional-111}
Under the hypotheses of~\cref{cor:conditional-112}, assume additionally that
the separately certified color-neutral first-phase reduction is compatible
with the same 114-move chain and saves one further move.  Then every Megaminx
state is solvable in at most 111 FTM moves.
\end{corollary}

\begin{proof}
Apply the additional one-move saving to the conditional bound in
\cref{cor:conditional-112}.
\end{proof}

The two local theorems are certified by this artifact.  The corollaries are
deliberately conditional because this work neither regenerates the complete
external 114-move computation nor proves compatibility of every surrounding
phase transition.
